\section{Positive transport and cyclic descent}

Let $t$ be the coordinate on a complex disk, oriented by increasing
argument. A positive meridian is the loop $t=\epsilon e^{2\pi iv}$,
$0\le v\le1$, with $\epsilon>0$. The integral monodromy of a family
along this loop is the map on the homology of the starting fibre
obtained by transporting cycles and identifying the ending fibre
with the starting fibre. Fix the cyclic change of coordinates $t=s^3$
and the generator $s\mapsto\zeta s$, where
$\zeta=e^{2\pi i/3}$. A deck transformation is an automorphism of
the pulled-back family covering this rotation and acting trivially
on its quotient.

\begin{lemma}\label{lem:inverse-descent}
Suppose a smooth family over the $s$ disk has a deck transformation
$\rho$ of order three. Mark its integral first homology by parallel
continuation over the disk, and let $A$ be the matrix of $\rho$ in
that marking. The positive monodromy of the quotient over $t\ne0$
is $A^{-1}$.
\end{lemma}
\begin{proof}
The positive meridian lifts from $s=\epsilon^{1/3}$ to
$s=\zeta\epsilon^{1/3}$. A homology class remains constant in the
disk marking along this lifted path. To express its endpoint in the
starting fibre of the quotient one applies $\rho^{-1}$. This acts
by $A^{-1}$. A translation of a torus is homotopic to the identity
through translations, so the same argument applies to an affine
deck transformation on its first homology. It makes no assertion
that the affine translation itself vanishes.
\end{proof}

The inverse can be seen without a monodromy table. Consider the
projective elliptic family whose affine equation is
\begin{equation}\label{eq:weierstrass}
 y^2=x^3-3t^3(t-1)x+2t^4(t-1)^2.
\end{equation}
Its origin is the point at infinity. For $s\ne0$, put
$t=s^3$, $x=s^4X$, and $y=s^6Y$. Dividing by $s^{12}$ gives
\begin{equation}\label{eq:good-equation}
 Y^2=X^3-3s(s^3-1)X+2(s^3-1)^2.
\end{equation}
For an equation $Y^2=X^3+aX+b$ its discriminant is
$-16(4a^3+27b^2)$; it vanishes exactly when the cubic and its
derivative have a common root. Here it is
$1728(s^3-1)^3$, so the projective fibres are smooth on a
sufficiently small disk around zero. The central curve is
$Y^2=X^3+2$. The deck transformation that fixes $x,y,t$ is
\begin{equation}\label{eq:actual-elliptic-deck}
 (X,Y,s)\longmapsto(\zeta^2X,Y,\zeta s).
\end{equation}
Substitution in \eqref{eq:good-equation} verifies the transformation
on the whole smooth family, including the origin section.

On the central curve its pullback on the holomorphic differential
$dX/Y$ is multiplication by $\zeta^2$. Thus for any cycle $c$,
\[
 \int_{\rho_*c}\frac{dX}{Y}=\zeta^2\int_c\frac{dX}{Y}.
\]
In a period marking with basis $(1,\zeta)$ the central deck matrix
is consequently
$\left(\begin{smallmatrix}-1&1\\-1&0\end{smallmatrix}\right)$.
The positive return matrix is its inverse. This calculation fixes
both the disk orientation and the direction of action on periods.

For the four-torus, retain the following explicit marked lattice
in $\mathbb C^2$. Choose $\tau$ with $\operatorname{Im}\tau>0$ and set
\begin{equation}\label{eq:period-lattice}
 e_1=(1,0),\quad e_2=(\zeta,0),\quad
 e_3=\tfrac13(2+\zeta,1),\quad e_4=(0,\tau),
 \qquad \Lambda=\sum_{j=1}^4\mathbb Ze_j.
\end{equation}
The vector $\delta=(0,1)=-2e_1-e_2+3e_3$ is an integral
period. The map multiplying the first coordinate by $\zeta$ and
fixing the second preserves $\Lambda$ and has matrix
\begin{equation}\label{eq:positive-T}
 T=\begin{pmatrix}0&-1&-1&0\\1&-1&0&0\\0&0&1&0\\0&0&0&1\end{pmatrix}.
\end{equation}
Indeed its values on the four basis vectors are
$e_2,-e_1-e_2,e_3-e_1,e_4$. Its inverse is the integral map
\begin{equation}\label{eq:inverse-T}
 A=T^{-1}=
 \begin{pmatrix}-1&1&-1&0\\-1&0&-1&0\\0&0&1&0\\0&0&0&1\end{pmatrix},
\end{equation}
with values $-e_1-e_2,e_1,e_3-e_1-e_2,e_4$.
In particular it fixes $\delta$ and $e_4$.

If \eqref{eq:positive-T} is the prescribed positive monodromy in
this unchanged disk marking, Lemma~\ref{lem:inverse-descent}
forces \eqref{eq:inverse-T} as the linear central deck action.
The elliptic calculation \eqref{eq:actual-elliptic-deck} agrees
with its action on the first complex coordinate. Conversely, the
elliptic calculation alone does not construct the lift to a
four-torus or identify an affine origin there. The distinction
between the linear map and an affine lift will matter below.

We first compute the complete chains for the quotient with deck
matrix $T$. We then use the inverse deck matrix to obtain positive
transport $T$ with the same fibre marking. Both quotients are
explicit geometric spaces; their identification with any prescribed
degeneration requires the corresponding affine descent data.
